% !TEX TS-program = lualatex+se
% !TEX encoding = UTF-8

\documentclass[oneside, 12pt, a4paper]{book}
%\usepackage[letterpaper,showframe]{geometry}
\usepackage{fullpage} % redukcja marginesów
\usepackage{geometry}
\geometry{
	left=15mm,
	right=10mm,
	top=12mm,
	bottom=15mm,
	marginparwidth=0mm,
	marginparsep=0mm
	}
\usepackage{gregoriotex}
\gresetlinecolor{gregoriocolor}
\gresetheadercapture{annotation}{greannotation}{string}
% Format annotation above initial
\grechangestyle{annotation}{\small\bfseries}%
\grechangedim{intersyllablespacenotes}{0.18cm}{scalable}%

\usepackage{fullpage}
\usepackage{multicol}

%  Konfiguracja kolumn
\usepackage{paracol}
\setcolumnwidth{13cm,5cm}
\setlength{\columnsep}{0.5cm}
\footnotelayout{m}
% \setlength{\columnseprule}{0.4pt}
% \colseprulecolor{red}
\usepackage{tcolorbox}
\newtcolorbox{nutybox}{size=minimal, colback=white, colframe=red, arc=0mm, width=\linewidth, halign=left, valign=center}
\newtcolorbox{tlumbox}{size=minimal, colback=white, colframe=red, arc=0mm, width=\linewidth, halign=left, valign=top, fontupper=\linespread{.8}\selectfont}

\usepackage{Tabbing}
\usepackage{longtable}
\usepackage[savepos]{zref}

\let\grelocalleftbox\localleftbox
\let\grelocalrightbox\localrightbox

\usepackage[latin,polish]{babel}

\usepackage{fontspec}
\defaultfontfeatures{Ligatures=TeX}
\setmainfont{Libertinus Serif}

\pagestyle{plain}	% usuwanie nagłówków/stopek z\,dokumentu
\usepackage[indent=0pt]{parskip} % ustawianie paragrafów
\usepackage{setspace} % odstępy między liniami
%\setstretch{1.5} % półtorej linii odległości między wierszami

\usepackage{graphicx} % umożliwia import obrazów
\graphicspath{{img/}} % katalog obrazów
\usepackage{pgfornament}% ornamenty graficzne

\usepackage{titlesec} % ustawianie nagłówków

\newfontface\vectis{VectisMonumentalW01-Regular}
\def\greinitialformat#1{\fontsize{20}{20}\vectis #1}

% --- komendy własne ---
\newcommand{\red}[1]{\textcolor{red}{#1}}
\newcommand{\black}[1]{\textcolor{black}{#1}}

\newcommand{\koniec}{
	\begin{center}
%		\vspace{5mm}
%		\pgfornament[width=0.6cm]{#1}
		\newpage
	\end{center}}

\begin{document}
\titlespacing*{\chapter}{0pt}{0pt}{5pt}	
\titleformat{\chapter}[hang]{\filcenter\LARGE\scshape\bfseries}{}{0pt}{} % format nagłówków
\titlespacing*{\section}{0pt}{0pt}{0pt}	
\titleformat{\section}[hang]{\filcenter\large\scshape}{}{0pt}{} % format nagłówków

\setlength{\parindent}{0pt}
\gresetbracerendering{font}

\pagenumbering{gobble} % nie jest wyprowadzany żaden numer strony, chociaż numer jest nadal resetowany. Odniesienia do tej strony również są puste.
\large
\begin{spacing}{0.7}
\begin{center}
	\section*{Fortem virili pectore (Commune non Virginum)}
	\small\textit{Liber Usualis, 1961, p.1234 \& Liber antiphonarius, 1960, p.99}
\end{center}	
\end{spacing}
\vspace{15mm}
\begin{paracol}{2}
	\ensurevspace{4cm}
	\begin{nutybox}\greannotation{Hymn.} \greannotation{2.} \gabcsnippet{(f3) FOr(f)tem(h) vi(hg)rí(i)li(hf) pé(g')cto(i)re(h'_) (,) Lau(h)dé(f')mus(g) o(h)mnes(gh) fé(f')mi(g)nam,(f.) (;) Quae(f) san(h)cti(hg)tá(i)tis(ji) gló(g')ri(i)a(h'_) (,) U(h)bí(f')que(g) ful(h)get(gh) ín(f')cly(e)ta.(f.) (::)}
	\end{nutybox}
	\switchcolumn
	\begin{tlumbox}
	\vspace{3.5mm} 
	Niewiastę mężną serca odważnego wysławiajmy wszyscy, która chwałą świętości jaśnieje wszędzie sławna.
	\end{tlumbox}
  	\ensurevspace{4cm}
	\switchcolumn*
	\begin{nutybox}
	\gresetinitiallines{0}\gabcsnippet{(c4)  <c>2.</c> Haec(f) san(h)ct{<e>o</e>}(h) a(hg)mó(i)re(hf) sáu(g')ci(i)a,(h'_) (,) Dum(h) mun(f')d{<e>i</e>}(g) a(g)mó(h)rem(gh) nó(f')xi(g)um(f.) (;) Hor(f)ré(h)scit,(hg) ad(i) cae(ji)lé(g')sti(i)a(h'_) (,) I(h)ter(f') per(g)é(h)git(gh) ár(f')du(e)um.(f.) (::)}
	\end{nutybox}
	\switchcolumn
	\begin{tlumbox}
	\vspace{3.5mm} 
	Ona (to) zraniona miłością świętą, gdy wzdrygała się przed zgubną miłością świata, przeszła stromą drogę do niebios.
	\end{tlumbox}
	\ensurevspace{4cm}
	\switchcolumn*
	\begin{nutybox}
	\gresetinitiallines{0}\gabcsnippet{ (c4) <c>3.</c> Car(f)nem(h) do(hg)mans(i) je(hf)jú(g')ni(i)is,(h'_) (,) Dul(h)cí(f')que(g) men(h)tem(gh) pá(f')bu(g)lo(f.) (;) O(f)ra(h)ti(hg)ó(i)nis(ji) nú(g')tri(i)ens,(h'_) (,) Cae(h)li(f') po(g)tí(h)tur(gh) gáu(f')di(e)is.(f.) (::)}
	\end{nutybox}
	\switchcolumn
	\begin{tlumbox}
	\vspace{3.5mm} 
	Ciało poskramiając postami i\,ducha karmiąc słodkim pokarmem modlitwy, bierze w posiadanie radości niebieskie.
	\end{tlumbox}
	\ensurevspace{4cm}
	\switchcolumn*
	\begin{nutybox}
	\gresetinitiallines{0}\gabcsnippet{(c4) <c>4.</c> Rex(f) Chri(h)ste,(hg) vir(i)tus(hf) fór(g')ti(i)um,(h'_) (,) Qui(h) ma(f')gna(g) so(h)lus(gh) éf(f')fi(g)cis,(f.) (;) Hu(f)jus(h) pre(hg)cá(i)tu,(ji) quaé(g')su(i)mus,(h'_) (,) Au(h)di(f') be(g)ní(h)gnus(gh) súp(f')pli(e)ces.(f.) (::)}
	\end{nutybox}
	\switchcolumn
	\begin{tlumbox}
	\vspace{3.5mm}
	Królu Chryste, dzielności mężnych który sam sprawiasz (rzeczy) wielkie, za jej orędownictwem, prosimy, wysłuchaj łaskawie korzących się. 
	\end{tlumbox}
	\ensurevspace{4cm}
	\switchcolumn*
	\begin{nutybox}
	\gresetinitiallines{0}\gabcsnippet{(c4) <c>5.</c> De(f)o(h) Pa(hg)tri(i) sit(hf) gló(g')ri(i)a,(h'_) (,) E(h)jús(f')que(g) so(h)li(gh) Fí(f')li(g)o,(f.) (;) Cum(f) Spí(h)ri(hg)tu(i) Pa(ji)rá(g')cli(i)to,(h'_) (,) Nunc(h) et(f') per(g) o(h)mne(gh) saé(f')cu(e)lum.(f.) (::) 
A(fgf)men.(ef..) (::)}
	\end{nutybox}
	\switchcolumn
	\begin{tlumbox}
	\vspace{3.5mm} 
	Bogu Ojcu niechaj będzie chwała i\,jedynemu jego Synowi wraz z\,Duchem Pocieszycielem teraz i\,po wszystek wiek. Amen.\footnotemark
	\end{tlumbox}
	\footnotetext{tłum. Ks. Bronisław Gładysz, "Hymny brewiarza rzymskiego oraz patronału polskiego. Przekład i objaśnienia", Poznań 1933} 
\end{paracol}	
\koniec		
\end{document}